%----------------------------------------------------------------------
%% APP: STRICT
%----------------------------------------------------------------------
% !TEX root = ./Main.tex


%We will fix an \emph{arbitrary} state and thus we will boil down to difference of scores for each player of any other strategy and the Nash policy for the specific state.
%We need to prove the following:

%\begin{itemize}
%    \item Difference in scores implies neighborhood in the primal space
%    \item Define $W_{init}$.
%    \item Write score differences and control the different terms (as before).
%    \item Select $M_{init}$ and the proof is complete.
%    \item Show difference of scores $\to -\infty$ convergence.
%\end{itemize}

%%%%%%Start here erase everything%%%
%In this section we will prove faster local convergence rates for deterministic Nash policies, when agents employ \cref{eq:LPG}.
%In the results that follow, we will fix an \emph{arbitrary} state $\rstate$ and we will focus on the update of the policies in this arbitrary state. Furthermore, we will drop the index $\rstate$ for convenience of notation.
%\begin{remark}
%In what follows when we write $\dstate_{\pure(\rstate),\run}$, we mean the  aggregated gradient that corresponds to the pure strategy profile of all agents, in the state $\rstate$ for the round $\run$. When we fix a state $\rstate$,  we instantly omit it and it is implied everywhere.
%\end{remark}


%----------------------------------------------------------------------
%% Preliminaries
%----------------------------------------------------------------------
\subsection{Structural preliminaries}

To prove \cref{thm:strict}, we will first require some notions describing the geometry of $\policies$ near $\eq$.
Referring to \cite{RW98} for a full treatment, we have:

\begin{definition}
Let $\cvx$ be a convex set and let $x\in \cvx$.
Then the tangent cone $\tcone_{\cvx}(x)$ is defined as the set of all rays emanating from $x$ and intersecting $\cvx$ to at least one other point different from $x$. The \textit{polar cone} $\pcone_{\cvx}(x)$ to $\cvx$ at $\point$ is then defined $\pcone_{\cvx}(x) = \braces{y: \braket{\dpoint}{\tvec} \leq 0 \text{ for all } \tvec\in \tcone_{\cvx}(x)}$, where $\dpoint$ belong in the dual space of the vector space in which $\cvx$ is defined.
\end{definition}

With these general definitions in hand, we proceed to characterize some further projections of Euclidean projections on $\policies$ that will play an important role in the sequel.
For notational simplicity, we suppress the player and state indices in the statement and proof of the next lemma.

\begin{lemma}
\label{lem:KKT}
$\point=\proj(\dpoint)$ if and only if there exist $\mu\in \R$ and $\nu_\pure\in\R_{+}$ such that, for all $\pure\in \pures$, we have
$\dpoint_\pure= \point_\pure+\mu-\nu_\pure$
with $\nu_\pure \geq 0$ and $\point_\pure \nu_\pure=0$.
\end{lemma}

\begin{proof}
Recall that 
\(
\proj(\dpoint)
    = \argmin_{\point\in\simplex(\pures)} \norm{\dpoint-\point}^{2}.
\)
Our result then follows by applying the KKT conditions to this optimization problem and noting that, since the constraints are affine, the KKT conditions are sufficient for optimality.
Our Langragian is
\begin{equation}
\mathcal{L}(\point,\mu,\nu)= \sum_{\pure\in\pures}\frac{1}{2}(\dpoint_\pure - \point_\pure )^2-{\mu(\sum_{\pure\in\pures} \point_\pure-1)} + {\sum_{\pure\in\pures} \nu_\pure \point_\pure}
\end{equation}
where the set of constraints (i) of the statement of the lemma  are the stationarity constraints, 
which in our case are $\nabla\mathcal{L}(\point,\mu,\nu)=0\Leftrightarrow
\nabla( \sum_{\pure\in\pures}\frac{1}{2}(\dpoint_\pure - \point_\pure )^2 )=\mu\nabla(\sum_{\pure\in\pures} \point_\pure-1) - \sum_{\pure\in\pures} \nu_\pure\nabla \point_\pure$
, while the set of constraints (ii) of the statement of the lemmas are the complementary slackness constraints.
Note that complementary slackness implies $\nu_\pure>0$ whenever $\pure\notin \supp(\point)$, so our proof is complete.
\end{proof}

Our next result is a concrete consequence of \cref{prop:var} which will be very useful in establishing the stability estimates required for the proof of \cref{thm:strict}.

\begin{lemma}\label{lem:basin of drift}
Let $\np = (\pure^*_{\play,\rstate})_{\play\in\players, \rstate\in\states}$ be a strict Nash policy.
Then there exists a neighborhood $\nhd$ of $\np$ and constants $c_{\play,\rstate}$ such that for each player $\play\in\players$ and state $\rstate\in\states$, we have:
\begin{equation}\label{eq:variational-inequality}
\text{$\payv_{\play\pure_{\play,\rstate}^*}\parens{\policy} - \payv_{\play\pure_{\play,\rstate}}\parens{\policy} \geq c_{\play,\rstate}\;$ for all $\policy\in\nhd$ and $\pure_\play\neq\pure_\play^*$, $\pure_\play\in\pures_\play$.}
\end{equation} 
\end{lemma}


\begin{proof}
Our claim  is a consequence of the definition of strict Nash policies. Specifically, from \cref{prop:var} we have
% \eqref{eq:sharp} we have that $\gvalue(\np)\in ri(PC(\np))$\footnote{with $ri$ we denote the relative interior.}, which means that 
\begin{equation}
    \inner{\gvalue(\np)}{z} <0 \quad\text{ for all }\quad z\in \tcone(\np), z\neq 0
\end{equation}
Let $z = e_{\play,\pure_{\play,\rstate}} -e_{\play,\pure^*_{\play,\rstate}}$, then we get that 
\begin{equation}
    \payv_{\play\pure_{\play,\rstate}^*}\parens{\np} - \payv_{\play\pure_{\play,\rstate}}\parens{\np} >0
\end{equation}
where $e_{\play,\pure_{\play,\rstate}}$ is the vector that has one only in the index and zero anywhere else.
By continuity there exists a neighborhood $\nhd\subseteq\pspace$ and $c_{\play,\rstate} > 0$ for each player $\play\in\players$ such that 
\begin{equation*}
 \payv_{\play\pure_{\play,\rstate}^*}\parens{\policy} - \payv_{\play\pure_{\play,\rstate}}\parens{\policy} \geq c_{\play,\rstate} \quad\text{ for all }\quad\policy\in\nhd
 \qedhere
\end{equation*}
\end{proof}

Our final result is intimately tied to the lazy projection step in \eqref{eq:LPG}, and quantifies the relation between initializations in $\dspace$ and $\policies$. 

\begin{lemma}\label{lem:equivalence} Let $\np = (\pure^*_{\play,\rstate})_{\play\in\players, \rstate\in\states}$, be a  deterministic policy. For each agent $\play\in\players$ and each state $\rstate\in\states$, let $\dstate_{\play,\pure_{\play,\rstate}} -\dstate_{\play,\pure^*_{\play,\rstate}}$ be the difference of the aggregated gradients between the strategy of the equilibrium and any other strategy $\pure_\play^*\neq \pure_{\play}\in\pures_\play$. Then for any $\eps >0 $ such that $\nhd_{\eps} = \braces{\policy: \policy_{\play,\pure_{\play,\rstate}^*}\geq 1-\eps \text{ for all } \play\in\players \text{ and } \rstate\in\states}$, there exist $M_{\play,\eps,\rstate}$ such that if $\dbasin_{\play,\rstate} = \braces{\dstate\in\R^{\pures_\play}: \dstate_{\play,\pure_{\play,\rstate}} -\dstate_{\play,\pure^*_{\play,\rstate}} < - M_{\play,\eps,\rstate}}$  then $\prod_{\play\in\players,\rstate\in\states}\proj_{\policies_{\play}}\parens{\dbasin_{\play,\rstate}}\subseteq \nhd_{\eps}$.
\end{lemma}
\begin{proof}

Consider an arbitrary player $\play\in\players$, a state $\rstate\in\states$, and let $\dbasin_\play\parens{M_{\play,\eps,\rstate}}$ be an open set as defined in the statement of the lemma. For notational simplicity, we will drop the index $\rstate$.
We will show that any $M_{\play,\eps}> 1-\frac{\eps}{|\pures_\play|}>0 $ satisfies our claim.
By  using \cref{lem:KKT} for a $\dstate_\play\in \dbasin_\play\parens{M_{\play,\eps}}$ with $\policy_\play=\proj\parens{\dstate_\play}$ we have that 
\begin{equation}
\begin{aligned}
    &\dstate_{\play\pure_\play^*}-\dstate_{\play\pure_\play}>M_{\play,\eps}\\
    &  {\policy_{\play\pure_\play^*}}-{\policy_{\play\pure_\play}} - (\nu_{\pure^*_\play}-\nu_{\pure_\play}) > M_{\play,\eps}\label{eq:contradiction}
\end{aligned}
\end{equation}
with $\nu_{\pure_\play} \geq 0$ and $\policy_{\play\pure_\play} =0$ whenever $\nu_{\pure_\play} >0$.
Notice that since $M_{\play,\eps} >1-\frac{\eps}{\nPures_{\play}}$  we  have that ${\policy_{\play\pure_\play^*}}>{\policy_{\play\pure_\play}} + 1 - \frac{\eps}{\nPures_{\play}} 
+ (\nu_{\pure^*_\play}-\nu_{\pure_\play})$ or 
\begin{equation}
  {\policy_{\play\pure_\play}}<   {\policy_{\play\pure_\play^*}} - 1 + \frac{\eps}{\nPures_{\play}} 
- (\nu_{\pure^*_\play}-\nu_{\pure_\play}) <\frac{\eps}{\nPures_{\play}} 
\end{equation}
Hence, by summing over all strategies of player $\play$ we get the desired result.
\end{proof}


%----------------------------------------------------------------------
%% Preliminaries
%----------------------------------------------------------------------
\subsection{Proof of the main theorem}

We are now in a position to prove our main result on the rate of convergence towards strict Nash policies.
For ease of reference, we restate \cref{thm:strict} below.

\strict*

\begin{proof}[Proof of \cref{thm:strict}]

We start by fixing a confidence level $\conf>0$ and all the parameters of the algorithm, such that all the assumptions stated in the theorem are satisfied and. We will prove that for each agent $\play\in\players$, $\rstate\in\states$ there exist $M_{1,\play,\rstate}>0$, $\dbasin_{1,\play,\rstate} = \braces{\dstate\in \R^{\pures_\play}: \dstate_{\play,\pure_\play} -\dstate_{\play,\pure^{*}_{\play}}<-M_{1,\play,\rstate} \text{ for all } \pure_\play\in\pures_\play, \pure_\play\neq\pure_{\play}^*}$,  such that if $\dstate_{\start}\in\dbasin_1 \defeq \prod_{\play\in\players,\rstate\in\states}\dbasin_{1,\play,\rstate}$  then the agents' sequence of play, converge to the deterministic Nash policy, in finite number of iterations. 

To simplify the notation, we will drop the indices $\rstate$ and $\play$ referring to the states and agents, accordingly, and we will focus on a specific agent and a specific state.
From \cref{lem:equivalence}, \cref{lem:basin of drift} we have that there exist constants $c,M$, neighborhood $\nhd_c =\braces{\policy\in\policies: \norm{\policy-\np}\leq \rad}$ and open set $\dbasin_{M}$ such that \begin{equation}
\begin{alignedat}{3}
      \payv_{\pure^*}\parens{\policy} -\payv_{\pure}\parens{\policy} &\geq c &\hspace{0.5em}&\text{ for all } \pure\neq\pure^*, \pure\in\pures \text{ and }\policy\in\nhd_c\\
      \dstate_{\pure^*}-\dstate_\pure &> M_c&\hspace{0.5em} &\text{ for all } \pure\neq\pure^*, \pure\in\pures  \text{ and } \policy = \proj\parens{\dstate}\in\nhd_c
\end{alignedat}
\end{equation}
The first step is to prove that for an appropriate initialization for $\dstate_\start$, we have $\dstate_\run\in \dbasin(M_c)$ for all $\run = \running$, with probability at least $1-\conf$. Assume that $\dstate_\runalt\in\dbasin(M_c)$ for all $\runalt=1,\hdots,n$; then 
for   the differences of the scores at a round $\run+1$ between any $\pure\in\pures$ and the equilibrium strategy $\pure^*$, we have
\begin{align}\label{eq: basic equality}
\dstate_{\pure,\run +1} -\dstate_{\pure^*,\run+1}
    &= \dstate_{\pure,\run} -\dstate_{\pure^*,\run} + (\signal_{\pure,\run} -\signal_{\pure^*,\run})
    \notag\\
    &=\dstate_{\pure,\start} -\dstate_{\pure^*,\start} + \sum_{\runalt = \start}^\run \step_\runalt [(\gvalue_{\pure,\runalt} -\gvalue_{\pure^*,\runalt})  + (\noise_{\pure,\runalt} -\noise_{\pure^*,\runalt}) +( \bias_{\pure,\runalt} -\bias_{\pure^*,\runalt})]
    \notag\\
    &\leq -M_1 + \sum_{\runalt = \start}^\run \step_\runalt [(\gvalue_{\pure,\runalt} -\gvalue_{\pure^*,\runalt})  + (\noise_{\pure,\runalt} -\noise_{\pure^*,\runalt}) +( \bias_{\pure,\runalt} -\bias_{\pure^*,\runalt})]
    \notag\\
    &\leq -M_1 -c\sum_{\runalt=1}^{\run}\step_{\runalt} + \sum_{\runalt=1}^{\run}\step_{\runalt}[ (\noise_{\pure,\runalt} -\noise_{\pure^*,\runalt}) +( \bias_{\pure,\runalt} -\bias_{\pure^*,\runalt})]
    \notag\\
    &\leq -M_1 -c\sum_{\runalt=1}^{\run}\step_{\runalt} + \sum_{\runalt=1}^{\run}\step_{\runalt}[ \snoise_{\runalt} + \sbias_{\runalt}]
\end{align}
where $\snoise_{\runalt} =  (\noise_{\pure,\runalt} -\noise_{\pure^*,\runalt})$ and $\sbias_{\runalt} =  2\norm{\bias_{\runalt}}$. 
Now, similarly to the proofs of \cref{thm:stable,thm:SOS} we will proceed to control the aggregate error terms 
\begin{equation}
\label{eq:err-agg}
\curr[R]
	= \sum_{\runalt=\start}^{\run} \iter[\step] \iter[\snoise]
	\quad
	\text{and}
	\quad
\curr[\submart]
	= \sum_{\runalt=\start}^{\run} {\iter[\step] \iter[\sbias]}.
\end{equation}
Since $\exof{\curr[\snoise] \given \curr[\filter]} = 0$, we have $\exof{\curr[R] \given \curr[\filter]} = \prev[R]$, so $\curr[R]$ is a martingale;
likewise, $\exof{\curr[\submart] \given \curr[\filter]} \geq \prev[\submart]$, so $\curr[\submart]$ is a sub-martingale. Furthermore from \eqref{eq:errorbounds}  we have:
\begin{enumerate}[I.]
     \item \label{eq:I}     $\exof{\snoise^{2}_\run} \leq  \exof{\norm{\noise_{\run}}^{2}} \leq \exof{\exof{\dnorm{\noise_{\run}}^{2} \given \filter_{\run}}} \leq \sdev^{2}_\run$\vspace{3pt}
     \item \label{eq:II}
     $\exof{\sbias_\run} = 2 \ex\bracks*{\dnorm{\bias_{\run}}}\leq \exof{\exof{\dnorm{\bias_{\run}} \given \filter_{\run}}} \leq \bbound_\run$
 \end{enumerate}

 Moreover, for any $\eta_{1} > 0$, we get by Doob's Maximal Inequality:

 \begin{align}
     \prob\parens*{\sup_{1 \leq \runalt \leq \run}R_{\runalt} \geq \eta_{1}} \leq \frac{\exof{R_{\run}^{2}}}{\eta_1^{2}} \stackrel{(a)}{=} \frac{\sum_{\runalt = 1}^{\run} \step_{\runalt}^{2}\exof{\snoise_{\runalt}^{2}}}{\eta_1^{2}} & \stackrel{(\ref{eq:I})}{\leq}
      \frac{\sum_{\runalt = 1}^{\run} \step_{\runalt}^{2}\sdev^{2}_{\runalt}}{\eta_1^{2}}   
 \end{align}
 where $(a)$ comes from the fact that $\exof{\snoise_i\snoise_j} = 0$ for $i \neq j$.  Since $\step_\run =\gamma/\run^{\pexp}$, $\sdev_\run = \bigoh(\run^{\sexp})$ and  $\pexp-\sexp >1/2$, 
 there exists $\gamma_1$ sufficiently small such that if $\gamma \leq \gamma_1$ then
 \begin{align}\label{eq: condition 1}
    \sum_{\runalt = 1}^{\infty}\step_{\runalt}^{2}\sdev^{2}_{\runalt} < \frac{\delta \eta_1^{2}}{2}
 \end{align}
 and so we automatically get that 
 \begin{equation}\label{eq:mart-prob}
     \prob\parens*{\sup_{1 \leq \runalt \leq \run}R_{\runalt} \geq \eta_{1}} \leq \frac{\conf}{2}
 \end{equation}
% %Hence, for $\step_{\run} = \frac{\step}{(\run + \runoff)^\pexp}$ and $\sdev_{\run} = \bigoh(\run^\sexp)$ with $\pexp - \sexp > 1/2$, and for $\runoff$ large enough we get the result.

Furthermore, notice that  the term $\{S_{\run}\}_{\run \in \N}$ is a sub-martingale, since $\exof{\abs{S_{\run}} \given \curr[\filter]} < \infty$ and $\exof{S_{\run+1} \given \curr[\filter]} > S_{\run}$, for all $\run$. As before, using Doob's Maximal Inequality, we get for any $\eta_2 > 0$:
 \begin{align}
     \prob\parens*{\sup_{1 \leq \runalt \leq \run}S_{\runalt} \geq \eta_{2}} \leq \frac{\exof{S_{\run}}}{\eta_{2}} 
     & = \frac{\sum_{\runalt = 1}^{\run} \step_{\runalt}\exof{\sbias_{\runalt}} }{\eta_2}
     {\leq}\frac{2\sum_{\runalt = 1}^{\run} \step_{\runalt}\bbound_{\runalt} }{\eta_2}
 \end{align}
So, since $\pexp+\bexp > 1$  there exists $\gamma_2$ sufficiently small such that if $\gamma \leq \gamma_2$ then
%\begin{equation}
\(
\sum_{\runalt = 1}^{\run} \step_{\runalt}\bbound_{\runalt} \leq \frac{\eta_2\conf}{4}
\)
%\end{equation}
which immediately implies that 
\begin{equation}\label{eq:prob-bias}
    \prob\parens*{\sup_{1 \leq \runalt \leq \run}S_{\runalt} \geq \eta_{2}} \leq\dfrac{\conf}{2}
\end{equation}
By choosing $\gamma \leq \min\braces{\gamma_1,\gamma_2}$ we get that 
 \begin{equation}
     \prob\parens*{\sup_{1 \leq \runalt \leq \run} R_\run + S_\run \leq M_c} \geq 1-\conf.
 \end{equation}
 Notice now that by choosing $M_1 > M_c + \eta_1 +\eta_2$, from \eqref{eq: basic equality} we have that with probability at least $1-\conf$, $\dstate_{\pure,\run+1} -\dstate_{\pure^*,\run+1} < -M_c$, which implies that $\policy_{\run+1}\in \nhd_c$.
 
 Defining the sequences of ``good'' events $\{\event_\run\}_{\run \in \N}$ and $\{\event'_\run\}_{\run \in \N}$ as $\event_\run \defeq \braces{\policy_\runalt \in \nhd_c \text{ for all }  \runalt=\running,\run}$ and $\event'_\run \defeq \braces*{\sup_{1 \leq \runalt \leq \run} R_{\runalt}+S_{\runalt} \leq \eta_1 + \eta_2}$, accordingly, we get that $\event'_\run \subseteq \event_\run$ for all $\run$. Because $\prob\parens*{\event'_\run} \geq 1-\conf$, we get that
% \begin{equation}
\(
\prob\parens*{\event_{\run}} \geq 1-\conf
\)
%\end{equation}
and since $\{\event_\run\}_{\run \in \N}$ is a decreasing sequence converging to $\event \defeq \braces*{\policy_\run \in \nhd_c, \forall \run \in \N}$, we obtain
%\begin{equation}
\(
\prob\parens*{\event} \geq 1-\conf.
\)
%\end{equation}
\ie
\begin{equation}
\prob\parens*{\curr \in \nhd_c, \; \forall \run \mid \dstate_{1} \in \dbasin_{1}} \geq 1-\conf
\end{equation}
Notice that the above conclusions immediately imply convergence in finite time. More specifically, constrained to the event $\event$ with probability at least $1-\conf$, from \cref{eq: basic equality} we have 
\begin{align}
    \dstate_{\pure,\run +1} -\dstate_{\pure^*,\run+1} \leq 
-M_c -c\sum_{\runalt=1}^{\run}\step_{\runalt} 
\end{align}
for all $\run=\running$.  Assume ad absurdum  that there exists at least one strategy $\pure\neq\pure^*,\pure\in\pures$ such that $\limsup_{\run\to\infty}\policy_{\pure,\run} \geq \eps >0$.
for all sufficiently large $\run$. Recall also that for $\policy\in \nhd_c$, it holds that $\policy_{\pure^*}>0$ by construction.
Using \cref{lem:KKT} we get
\begin{align}
   \dstate_{\pure,\run +1} -\dstate_{\pure^*,\run+1}=  \policy_{\pure,\run+1} - \policy_{\pure^*,\run+1}  &\leq -M_c -c\sum_{\runalt=1}^{\run}\step_{\runalt}
\end{align}
Notice that the \acs{LHS} of this inequality is bounded, while the \acs{RHS} goes to $-\infty$, which is a contradiction. Thus, with probability at least $1-\conf$, $\policy_{\run}\to\policy^*$ as $\run\to\infty$.

We can rewrite the previous inequality as
\begin{equation}
    \policy_{\pure,\run+1} \leq 1 -M_c -c\sum_{\runalt=1}^{\run}\step_{\runalt} \quad\text{ for all }\quad\pure^*\neq\pure\in\pures
\end{equation}
Now aggregating over all strategies, on the previous inequality,  we get that 
\begin{align}
    \onenorm{\policy_{\run+1} -\np} = 2(1-\policy_{\pure^*,\run+1}) \leq2 \sum_{\pure^*\neq\pure\in\pures}(1- M_c -c\sum_{\runalt=1}^{\run}\step_{\runalt} )
\end{align}
Thus, once $\sum_{\runalt=1}^{\run}\step_{\runalt}$  becomes at least $(1-M_c)/c$, which occurs in finite time, the convergence is implied.
\end{proof}


\begin{proof}[Proof of \cref{cor:strict}]
For \cref{mod:full,mod:stoch}, taking $\bexp = \infty, \sexp = 0$ we readily get that $\pexp > 1/2$. Since we require that $\sum_{\run=1}^{\infty}\step_\run = \infty$, we obtain that $\pexp \in (1/2,1]$. 
%Of course the larger, the step-size the faster the convergence in this case.

For \cref{mod:value}, we have that $\curr[\bbound] = \bigoh(\curr[\expar])$ and $\curr[\sdev] = \bigoh(1/\sqrt{\curr[\expar]})$, \ie $\bexp = r$ and $\sexp = r/2$. Now, since $\pexp \leq 1$, $\pexp + \bexp >1$ and $\pexp - \sexp > 1/2$, we obtain that $\pexp \in (2/3,1]$. 
\end{proof}